\par\Needspace{6\baselineskip}\medskip
\begin{center}\hypertarget{native-226-article163}{}\pdfbookmark[2]{163.}{article-163}163.\end{center}
Methodus, quam inde ab art. 136 exposuimus, praecipue quidem determinationi
primae orbitae penitus adhuc incognitae accomodata est: attamen successu aeque
felici tunc quoque in usum vocatur, ubi de correctione orbitae proxime iam
cognitae per tres observationes quantumvis ab invicem distantes agitur.
In tali autem casu quaedam immutare conveniet. Scilicet quoties observationes
motum heliocentricum permagnum complectuntur, haud amplius licebit,
$\dfrac{\theta''}{\theta}$ atque $\theta\theta''$ tamquam valores approximatos
quantitatum $P$, $Q$ considerare: quin potius ex elementis proxime cognitis
valores multo magis exacti elici poterunt. Calculabuntur itaque levi calamo
per ista elementa pro tribus observationum temporibus loca heliocentrica in
orbita, unde designando anomalias veras per $v$, $v'$, $v''$, radios vectores
per $r$, $r'$, $r''$, semiparametrum per \SourceLeaf{227}$p$, prodibunt valores approximati sequentes:
\[
\begin{aligned}
P&=\frac{r\sin(v'-v)}{r''\sin(v''-v')}\mathpunct{\raisebox{.45ex}{,}}\\[.6em]
Q&=\frac{4r'^4\sin\frac12(v'-v)\sin\frac12(v''-v')}{p\cos\frac12(v''-v)}\mathpunct{\cdot}
\end{aligned}
\]
\noindent His itaque hypothesis prima superstruetur, paullulumque ad libitum immutatis
secunda et tertia: haud enim e re esset, $P'$ et $Q'$ hic pro novis valoribus
adoptare (uti supra fecimus), quum hos valores magis exactos evadere haud
amplius supponere liceat. Hac ratione omnes tres hypotheses commodissime
\emph{simul} absolvi poterunt: quarta dein secundum praecepta art. 120 formabitur.
Ceterum haud abnuemus, si quis unam alteramve decem methodorum in artt.
124---129 expositarum in tali casu si non magis tamen aeque fere expeditam
existimet, ideoque in usum vocare malit.
\par\bigskip\begin{center}\rule{.23\linewidth}{.35pt}\end{center}
